% Theorem boxes framed by TikZ through ntheorem's [framed] option (\newshadedtheorem with
% \theoremframecommand) and framed.sty,
% as in "An Infinite Descent into Pure Mathematics" (book/includes/theorems.tex):
% each environment's body is boxed, measured by \fb@sizeofframe (which draws the
% whole TikZ frame around a rule) and framed, split across pages when long.
% Written for the corpus (lane P6-INFDESC-PAGE, BOX-MEMO's sweep coverage).
\documentclass[10pt]{article}
\usepackage[a5paper,margin=18mm]{geometry}
\usepackage{amsmath,amssymb}
\usepackage{xcolor}
\usepackage{tikz}
\usepackage{framed}
\usepackage[thmmarks,framed]{ntheorem}

\definecolor{thmbgcol}{RGB}{238,244,252}
\definecolor{thmbdcol}{RGB}{40,90,170}
\definecolor{defbgcol}{RGB}{246,240,252}
\definecolor{defbdcol}{RGB}{120,60,160}
\definecolor{pfbdcol}{RGB}{150,150,150}

\tikzstyle{thmbox} = [rectangle, fill=thmbgcol, inner xsep=4pt]
\newcommand\thmbox[1]{%
  \noindent\begin{tikzpicture}%
  \node [thmbox] (box){%
    \begin{minipage}{\linewidth}%
      #1%
    \end{minipage}%
  };%
  \draw[color=thmbdcol, line width=1.5pt] (box.south west) -- (box.north west) ;
  \end{tikzpicture}}
\tikzstyle{defbox} = [rectangle, fill=defbgcol, inner xsep=4pt]
\newcommand\defbox[1]{%
  \noindent\begin{tikzpicture}%
  \node [defbox] (box){%
    \begin{minipage}{\linewidth}%
      #1%
    \end{minipage}%
  };%
  \draw[color=defbdcol, line width=1.5pt] (box.south west) -- (box.north west) ;
  \end{tikzpicture}}
\tikzstyle{pfbox} = [rectangle, inner xsep=4pt]
\newcommand\pfbox[1]{%
  \noindent\begin{tikzpicture}%
  \node [pfbox] (box){%
    \begin{minipage}{\linewidth}%
      #1%
    \end{minipage}%
  };%
  \draw[color=pfbdcol, line width=0.75pt] (box.south west) -- (box.north west) ;
  \end{tikzpicture}}

\theoremstyle{break}
\theorembodyfont{\normalfont}
\theoremheaderfont{\bfseries\color{thmbdcol}}
\theorempreskip{0pt}\theorempostskip{0pt}
\theoremsymbol{}
\let\theoremframecommand\thmbox
\newshadedtheorem{theorem}{Theorem}[section]
\newshadedtheorem{lemma}[theorem]{Lemma}
\theoremheaderfont{\bfseries\color{defbdcol}}
\let\theoremframecommand\defbox
\newshadedtheorem{definition}[theorem]{Definition}
\theoremheaderfont{\itshape}
\theoremsymbol{\ensuremath{\square}}
\let\theoremframecommand\pfbox
\theoremnumbering{arabic}
\newshadedtheorem{proof}{Proof}

\begin{document}

\section{Equivalence relations}

A relation on a set is a way of saying which elements are related. We begin
with the definition and then prove that equivalence relations and partitions of
a set are two descriptions of the same thing.

\begin{definition}
Let $X$ be a set. A relation $\sim$ on $X$ is an \emph{equivalence relation}
if it is reflexive ($x \sim x$ for all $x \in X$), symmetric ($x \sim y$ implies
$y \sim x$) and transitive ($x \sim y$ and $y \sim z$ imply $x \sim z$).
\end{definition}

Given an equivalence relation, each element determines the set of elements
related to it, its \emph{equivalence class} $[x] = \{ y \in X \mid x \sim y \}$.
The classes cover $X$, since $x \in [x]$, and two of them are either equal or
disjoint.

\begin{lemma}
Let $\sim$ be an equivalence relation on $X$ and let $x, y \in X$. Then
$[x] = [y]$ if and only if $x \sim y$, and if $[x] \ne [y]$ then
$[x] \cap [y] = \varnothing$.
\end{lemma}

\begin{proof}
Suppose $[x] = [y]$. Since $y \in [y] = [x]$ we have $x \sim y$. Conversely,
suppose $x \sim y$ and let $z \in [y]$, so that $y \sim z$; by transitivity
$x \sim z$, so $z \in [x]$. Hence $[y] \subseteq [x]$, and by symmetry
$[x] \subseteq [y]$. Finally, if $z \in [x] \cap [y]$ then $x \sim z$ and
$y \sim z$, so $x \sim y$ by symmetry and transitivity, and $[x] = [y]$.
\end{proof}

\begin{theorem}
Let $X$ be a set. For every equivalence relation $\sim$ on $X$ the set
$X/{\sim} = \{ [x] \mid x \in X \}$ is a partition of $X$, and every partition
of $X$ arises in this way from exactly one equivalence relation.
\end{theorem}

\begin{proof}
That $X/{\sim}$ is a partition is the previous lemma together with $x \in [x]$.
Conversely, let $\mathcal{U}$ be a partition of $X$ and define $x \approx y$ if
and only if $x$ and $y$ lie in the same block of $\mathcal{U}$.
\begin{itemize}
\item Each $x$ lies in some block, so $x \approx x$.
\item If $x$ and $y$ share a block then so do $y$ and $x$.
\item If $x, y$ share a block $U$ and $y, z$ share a block $V$, then
$y \in U \cap V$, and since blocks are disjoint or equal, $U = V$; so $x, z$
share a block.
\end{itemize}
The equivalence classes of $\approx$ are exactly the blocks of $\mathcal{U}$:
\begin{align*}
[x]_{\approx} &= \{ y \in X \mid x \approx y \} && \text{by definition} \\
&= \{ y \in X \mid y \in U_x \} && \text{where $U_x$ is the block of $x$} \\
&= U_x .
\end{align*}
If two equivalence relations have the same classes they relate the same pairs,
so the relation is unique. This completes the proof, which is long enough to
show how a framed box is split across a page break: the frame is measured again
for the first and the last part, and each part is drawn with its own picture.
We add a few more lines of reasoning so that the box needs more room than the
page has left. Suppose that $f : X \to Y$ is any function; then $x \sim y$ if
and only if $f(x) = f(y)$ defines an equivalence relation whose classes are the
non-empty fibres $f^{-1}[\{y\}]$, and the quotient map $q : X \to X/{\sim}$ is a
surjection through which $f$ factors uniquely as $f = \bar{f} \circ q$ with
$\bar{f}$ injective. In particular $|X/{\sim}| = |f[X]|$.
\end{proof}

\section{Modular arithmetic}

\begin{definition}
Fix an integer $n > 0$. Integers $a$ and $b$ are \emph{congruent modulo $n$},
written $a \equiv b \pmod{n}$, if $n$ divides $a - b$.
\end{definition}

\begin{theorem}
Congruence modulo $n$ is an equivalence relation on $\mathbb{Z}$ with exactly
$n$ classes, $[0], [1], \dots, [n-1]$.
\end{theorem}

\begin{proof}
Reflexivity and symmetry are clear since $n \mid 0$ and $n \mid -(a-b)$. If
$n \mid a-b$ and $n \mid b-c$ then $n \mid (a-b)+(b-c) = a-c$. By the division
theorem every $a$ is $qn + r$ with $0 \le r < n$, so $a \equiv r$, and two
distinct remainders are not congruent since their difference has absolute value
less than $n$.
\end{proof}

\begin{lemma}
If $a \equiv a'$ and $b \equiv b' \pmod{n}$ then $a + b \equiv a' + b'$ and
$ab \equiv a'b' \pmod{n}$.
\end{lemma}

\begin{proof}
We have $(a+b) - (a'+b') = (a-a') + (b-b')$ and
$ab - a'b' = a(b-b') + b'(a-a')$, both multiples of $n$.
\end{proof}

\end{document}
